% Weekly exercise sheet for Lecture 10; also included in the synthesis sheet.
\ifdefined\ExercisesSynthesis
\else
\documentclass[11pt]{article}
\usepackage[utf8]{inputenc}
\usepackage[T1]{fontenc}
\usepackage[english]{babel}
\usepackage{lmodern,microtype}
\usepackage[margin=0.9in]{geometry}
\usepackage{amsmath,amssymb,amsthm,mathtools}
\usepackage{enumitem,hyperref}
\newcommand{\E}{\mathbb E}
\newcommand{\Qbb}{\mathbb Q}
\newcommand{\dd}{\,\mathrm d}
\newcommand{\F}{\mathcal F}
\theoremstyle{definition}
\newtheorem{exercise}{Exercise}[section]
\newlist{exparts}{enumerate}{1}
\setlist[exparts,1]{label=\textbf{(\alph*)},itemsep=0.3em,topsep=0.4em,leftmargin=2em}
\newcommand{\lecturesummary}[1]{\begin{quote}\small\textbf{Focus.} #1\end{quote}}
\title{Stochastic Calculus and Option Pricing\\Lecture 10 Exercises}
\author{}
\date{}
\begin{document}
\maketitle
\vspace{-1.2em}
\setcounter{section}{9}
\fi

\section{Week 10 --- American Options and Optimal Stopping}
\lecturesummary{This question extends optimal stopping to the writer's hedge
and requires a policy argument rather than a binomial price computation.}


\begin{exercise}[A supermartingale gives a writer's hedge]
Work in a discrete-time, arbitrage-free market with money-market account
\(B_t\). Let \(G_t\) be the amount due if an American claim is exercised at
date \(t\). Suppose a self-financing portfolio has value \(X_t\) and
\(X_t/B_t\) is a nonnegative supermartingale under a risk-neutral measure.
\begin{exparts}
  \item Prove that if \(X_t\geq G_t\) at every exercise date, then the
  initial portfolio value upper-bounds the expected discounted payment for
  every admissible exercise rule.
  \item Explain why checking only the terminal payoff \(G_T\) is
  insufficient for an American claim.
  \item What additional condition would make this superhedge cheapest among
  all portfolios that cover every exercise date? Relate your answer to the
  least supermartingale that dominates the discounted exercise payoff.
\end{exparts}
\end{exercise}

\ifdefined\ExercisesSynthesis\else
\end{document}
\fi
